\section*{第 \ref{ch:g8:intensity-ammeter} 章 --- 电流强度与电流表}

\begin{solution}{exo:g8:intensity-ammeter:1}
电流强度就是单位时间内通过的电荷量，
好比河流每秒流过的水量。单位是安培，\unit{A}；
$\qty{250}{mA} = \qty{0.25}{A}$。
\end{solution}

\begin{solution}{exo:g8:intensity-ammeter:2}
\qty{0.02}{A}：指示灯。\qty{0.3}{A}：手电筒灯泡。
\qty{10}{A}：电水壶。\qty{100}{A}：发动机启动时。
\end{solution}

\begin{solution}{exo:g8:intensity-ammeter:3}
电流表只计量流经自身的电流，因此必须
串联在电路中——即串联。它对电流的阻碍几乎
可以忽略；若直接跨接在电池两端，它便成了
一个带表盘的短路。
\end{solution}

\begin{solution}{exo:g8:intensity-ammeter:4}
断开开关，断开回路；接入电流表，电流由\unit{A}接线柱流入，
由公共端流出；先选最大量程，
再逐级减小；闭合开关，读数，写上单位。
先用最大量程，是为了在电流意外偏大时，
由最坚固的挡位承受，而不致损坏精密的挡位。
\end{solution}

\begin{solution}{exo:g8:intensity-ammeter:5}
在串联电路中，回路上各处的电流强度相同。
这把早先的观察化作了数字：两只相同的串联灯泡，
无论位于何处，亮度都完全相同。
\end{solution}

\begin{solution}{exo:g8:intensity-ammeter:6}
\qty{0.45}{A}，\qty{0.45}{A}：同一回路，电流强度相同。
\end{solution}

\begin{solution}{exo:g8:intensity-ammeter:7}
不对。定律保证的是\emph{同一时刻沿回路各处}的电流值相同，
而不是回路重新接线后仍保持同一数值。第二只灯泡使
电流减弱——但各处同等程度地减弱：现在所有的电流表都
读数为\qty{0.17}{A}。
\end{solution}

\begin{solution}{exo:g8:intensity-ammeter:8}
电流表接反了——电流从公共端流入。
若是数字电流表，无妨：唯一的“代价”是一个负号
（指针式电流表则会猛撞到止挡）。
\end{solution}

\begin{solution}{exo:g8:intensity-ammeter:9}
$0.60 - 0.35 = \qty{0.25}{A}$——节点必须守恒：
流入多少，就流出多少。第二个节点处的电流表
读数为汇合后的\qty{0.60}{A}。
\end{solution}

\begin{solution}{exo:g8:intensity-ammeter:10}
并联时，几乎毫无阻碍的电流表相当于一条旁路：灯泡
被弃用，逐渐变暗直至熄灭，而电流表中却流过
一股几乎不受限制的强电流——相当于经过仪表的短路，
运气好的话，只会烧断它内部的保险丝。
这名学生亲手造出了安全一章里的“头号反派”。
\end{solution}

\begin{solution}{exo:g8:intensity-ammeter:11}
干路中的电流等于各支路电流之和：\qty{0.60}{A}。再并入
一条相同的支路，总电流变为\qty{0.90}{A}——
此时电池要供给三路满载电流，消耗最快。
\end{solution}

\begin{solution}{exo:g8:intensity-ammeter:12}
需要知道的电流有三个：干路电流$I$，以及各支路电流
（$I_1$、$I_2$）——共三个（节点之外的回路段重复这些数值）。
只需两次测量即可：任测两个，
由节点关系$I = I_1 + I_2$即可求出第三个。安全的操作顺序是：断开电路；将电流表接入
所选支路；选最大量程；闭合；读数；断开后再移动
电流表——而且在任何步骤中，电流表都不得同时连接
两条不同支路上的两点。
%
\end{solution}

\begin{solution}{pb:g8:intensity-ammeter:1}
\textbf{1.} 断开电路；在灯泡与电池之间断开回路；
接入电流表，电流由\unit{A}端流入，由公共端流出；
选最大量程，逐级减小；闭合电路；读数并写上单位。
\textbf{2.} \qty{0.28}{A}。
\textbf{3.} 仍为\qty{280}{mA}——同一回路，电流强度相同：
电流表无论在哪一侧，计量的都是同一股电流。
\textbf{4.} 回路上任意位置——两灯串联，无论位置如何，
两表读数都是\qty{280}{mA}。
\textbf{5.} 定律否定了这一预测：串联灯泡通过的电流
完全相同。差异在于灯泡本身——
构造不同，会把相同的电流转化为不同的亮度
（后面的章节将说明具体差异）。
\textbf{6.} 回路各点电流均为\qty{150}{mA}——A 之前、
B 之后、电池旁。
\textbf{7.} 灯泡 B 烧毁，被旁路弃置，回路
失去了一个阻碍者，因而通过的电流增大：电流表
读数上升——灯泡 A 也相应变亮。
\textbf{8.} “跨接在 B 两端时，旁路仍使灯泡 A 守在
回路中；跨接在电池两端，则什么都不剩——
而没有守卫的电流属于火灾一章，而不是
一次演示。”
\textbf{9.} $0.52 - 0.30 = \qty{0.22}{A}$。
\textbf{10.} $0.40 + 0.30 = \qty{0.70}{A}$——而电机支路
的保险丝（或断路器）正密切“监听”电机线路——
电机卡住时，正是保护装置发挥作用的时候。
\textbf{11.} 灯泡支路：\qty{0}{A}；电机支路：仍为其自身的
\qty{0.40}{A}，不受影响（并联各支路互不影响）；干路：
\qty{0.40}{A}。
\textbf{12.} 电池要为“慷慨”付出代价：两条正常的支路
共需两路满载电流，而干路的总电流
会使电池最快耗尽。
\textbf{13.} 例如：“电流表应串联在电路中，
绝不能跨接在电路两端。同一回路保证电流强度处处相同——
任何地方，任何时候都是如此。节点的电流如同账目，流入
多少，就流出多少，精确到毫安。”
\end{solution}
